\documentclass[11pt]{article}
\usepackage[a4paper,margin=25mm]{geometry}
\usepackage{amsmath,amssymb,mathtools}
\usepackage{fontspec,unicode-math}
\setmainfont{Latin Modern Roman}
\setsansfont{Latin Modern Sans}
\setmonofont{Latin Modern Mono}
\setmathfont{Latin Modern Math}
\usepackage{xurl,hyperref}
\hypersetup{colorlinks=true,urlcolor=blue,linkcolor=blue}
\newcommand{\sref}[2]{\hyperlink{source:#1:#2}{[S#2]}}
\newcommand{\cataloguescope}[1]{\paragraph{Mathematical question.}#1}
\begin{document}
\section{Quantum algorithm for Smallest Conjugate}\label{statement}

\subsection{Definitions and mathematical statement}

\paragraph{Shared notation.}$\mathbb N=\{1,2,\ldots\}$, and $\mathbb Z,\mathbb Q,\mathbb R,\mathbb C$ denote the integers, rationals, reals and complex numbers. Any other conventions are specified locally.


Let $c$ be an algebraic integer, specified by its monic irreducible polynomial in $\mathbb Z[X]$, and let $S\subset\mathbb Z[c]$ be a nonzero ideal specified by an integral basis in the coefficient basis $1,c,\ldots,c^{n-1}$, where $n=[\mathbb Q(c):\mathbb Q]$. For $v\in\mathbb Q(c)$ define
\[
\|v\|_{\rm conj}=\max_{\sigma:\mathbb Q(c)\hookrightarrow\mathbb C}|\sigma(v)|.
\]
The Smallest Conjugate problem is minimization over nonzero ideal elements, rather than selection of a conjugate of one given element. The nonzero convention is the shortest-vector convention used throughout the source and in its immediate ideal-lattice equivalences.
\paragraph{Statement recorded in the supplied source.}
Develop a quantum algorithm, polynomial in the bit length of the defining polynomial and ideal basis, which returns $0\ne v\in S$ minimizing $\|v\|_{\rm conj}$. \sref{statement}{2}

\subsection{Short English statement}
Find an efficient quantum algorithm for the ideal element whose largest conjugate magnitude is smallest.
\subsection{Sources}
\begin{description}
\item[\hypertarget{source:statement:1}{S1}] ULAM.AI and the UnsolvedMath Contributors, UnsolvedMath: A Curated Collection of Open Mathematics Problems, catalogue entry. CC BY 4.0. \url{https://huggingface.co/datasets/ulamai/UnsolvedMath}.
\item[\hypertarget{source:statement:2}{S2}] Daniele Micciancio, \emph{Ideal Lattices: cryptographic applications and open problems}, IPAM Number Theory and Cryptography: Open Problems workshop (9--13 October 2006), PDF pages 28--34, 41. \url{https://helper.ipam.ucla.edu/publications/SCWS1/SCWS1_6240.pdf}.
\end{description}
\label{end:statement}
\end{document}
